% Generated from project outputs. Do not edit by hand.
\begin{table}[!htbp]
\centering
\begin{threeparttable}
\caption{Including nonlabor by subsample: 2016-2024}
\label{tab:descriptive-2016-2024-subsamples-y3}
\small
\setlength{\tabcolsep}{4pt}
\renewcommand{\arraystretch}{1.15}
\begin{tabularx}{\linewidth}{@{}Xrrrrr@{}}
\toprule
Sample & N & Mean t & SD t & Mean t+1 & SD t+1 \\
\midrule
General & 552 & 2,913.6 & 1,786.2 & 3,111.6 & 1,795.1 \\
Men & 552 & 3,262.7 & 2,087.6 & 3,488.7 & 2,070.2 \\
Women & 552 & 2,339.7 & 1,568.3 & 2,528.4 & 1,649.8 \\
Urban & 552 & 3,159.5 & 1,834.8 & 3,338.9 & 1,815.6 \\
Rural & 547 & 2,391.7 & 1,635.4 & 2,634.6 & 1,817.8 \\
Indigenous & 547 & 2,371.6 & 1,690.4 & 2,643.2 & 1,765.7 \\
Non-Indigenous & 460 & 3,088.7 & 1,862.1 & 3,223.3 & 1,819.7 \\
Bottom 99 percent & 552 & 2,693.0 & 1,408.9 & 2,887.0 & 1,431.2 \\
Bottom 90 percent & 552 & 2,229.3 & 985.6 & 2,383.7 & 983.9 \\
Bottom 50 percent & 548 & 1,266.5 & 539.1 & 1,367.0 & 535.5 \\
\bottomrule
\end{tabularx}
\begin{tablenotes}[flushleft]
\footnotesize
\item Monthly income in quetzales. Unweighted descriptive statistics of survey-weighted cohort means. N counts cohort-transition rows, not individuals or independent cohorts. SD is dispersion across cohort means, not the standard error of a mean. Six annual transitions are pooled; the 2019-2021 gap is excluded. Both incomes must be observed. Subsamples overlap; bottom-income groups are cumulative.
\end{tablenotes}
\end{threeparttable}
\end{table}
